\chapter{Host preparation and PCIe driver}
\label{chap:host-driver}

\section{Initial checks}

Before installing the driver, check the operating system version, architecture and
kernel. On Linux, you can save this information in \code{evidence}:

\begin{lstlisting}[language=bash]
mkdir -p evidence
{
  cat /etc/os-release
  uname -r
  uname -m
  lspci -nn
} | tee evidence/host-inventory.txt
\end{lstlisting}

\code{uname -m} usually reports \code{x86\_64} on a PC and \code{aarch64} on a
Raspberry Pi 5 with a 64-bit operating system. The kernel version is needed to check
driver compatibility.

\section{Enable PCIe on Raspberry Pi 5}

This section applies only to Raspberry Pi 5. Compatible HAT+ devices can be detected
automatically, whereas Raspberry Pi specifies that a non-HAT+ PCIe adapter needs the
PCIe interface enabled in the configuration file \cite{rpi_pcie}.

On current Raspberry Pi OS, open:

\begin{lstlisting}[language=bash]
sudoedit /boot/firmware/config.txt
\end{lstlisting}

If that path does not exist, check \code{/boot/config.txt}. Add:

\begin{lstlisting}[language=bash]
dtparam=pciex1
\end{lstlisting}

Save the file and reboot:

\begin{lstlisting}[language=bash]
sudo reboot
\end{lstlisting}

Raspberry Pi 5 uses PCIe Gen2 by default. Gen3 can be enabled, but Raspberry Pi warns
that the board is not certified for stable operation at that speed \cite{rpi_pcie}.
Keep Gen2 for bring-up.

\section{Check the PCIe connection}

After booting, check that the system detects the board:

\begin{lstlisting}[language=bash]
lspci -nn
\end{lstlisting}

A PCIe device corresponding to the board should appear. If it does not, check the
assembly, power and, on Raspberry Pi 5, whether PCIe is enabled before proceeding
with the driver.

For more information about the detected device, run:

\begin{lstlisting}[language=bash]
lspci -nnk
\end{lstlisting}

This stage checks only the PCIe connection. Seeing the board in \code{lspci} does
not yet mean that the Akida driver is installed.

\section{Install the driver}

BrainChip's official PCIe driver is distributed in the \code{akida\_dw\_edma}
repository. Before installing it, check the kernel version:

\begin{lstlisting}[language=bash]
uname -r
\end{lstlisting}

The driver version consulted for this manual supports kernels from 5.4 to 6.8.
The module does not build on kernels 6.9 or later \cite{brainchip_driver}.
Compatibility can change, so also check the repository README before installation.

\begin{stopbox}{Unsupported kernel}
If the kernel version is outside the driver's supported range, do not proceed with
installation until you have checked whether BrainChip has released a compatible version.
\end{stopbox}

On Ubuntu, install the dependencies specified by BrainChip:

\begin{lstlisting}[language=bash]
sudo apt update
sudo apt install dkms build-essential "linux-headers-$(uname -r)"
\end{lstlisting}

Package names may differ on other distributions, but DKMS, build tools and headers
for the running kernel must be installed.

Then clone the repository and install the driver:

\begin{lstlisting}[language=bash]
git clone https://github.com/Brainchip-Inc/akida_dw_edma.git
cd akida_dw_edma
git rev-parse HEAD
./install.sh
\end{lstlisting}

The output of \code{git rev-parse HEAD} records the exact driver revision used;
you can also use your own configuration.

\section{Check the driver}

After installation, check the DKMS status and confirm that the Akida device is available:

\begin{lstlisting}[language=bash]
dkms status
lsmod | grep -i akida
ls -l /dev/akida*
lspci -nnk
\end{lstlisting}

Installation is successful when the device still appears on PCIe, the module is loaded
and at least one \code{/dev/akida*} node exists.

If the device appears in \code{lspci} but \code{/dev/akida*} is missing, read the
kernel messages:

\begin{lstlisting}[language=bash]
sudo dmesg | grep -iE 'akida|edma|pcie' | tail -n 120
\end{lstlisting}

If the node exists but is accessible only as \code{root}, check its permissions and
the udev rule installed by the driver before proceeding.

With the driver working, continue with the software environment setup in
\cref{chap:software}.
